\documentclass[10pt,twocolumn]{article}

% --- Packages ---
\usepackage[margin=1in]{geometry}
\usepackage{amsmath,amssymb}
\usepackage{booktabs}
\usepackage{hyperref}
\usepackage{microtype}
\usepackage{graphicx}
\usepackage{xcolor}
\usepackage{enumitem}
\usepackage{caption}
\usepackage{float}
\usepackage{cite}
\usepackage{url}
\usepackage{parskip}
\setlength{\columnsep}{0.25in}

\hypersetup{
  colorlinks=true,
  linkcolor=blue!70!black,
  citecolor=blue!70!black,
  urlcolor=blue!70!black
}

% --- Title ---
\title{\textbf{Gemmulem: A High-Performance Pure-C Library for\\
Finite Mixture Model Deconvolution via the EM Algorithm}}

\author{
  Micah Thornton\textsuperscript{1} \and Chanhee Park\textsuperscript{1}
  \\[4pt]
  \small\textsuperscript{1}Department of Statistics, [Institution]\\
  \small\texttt{\{mathornton01, chanheepark\}@[institution].edu}
}

\date{}

\begin{document}

\maketitle
\thispagestyle{empty}

% =============================================================================
\begin{abstract}
Finite mixture models are a cornerstone of unsupervised learning and density
estimation, yet existing high-quality implementations either carry heavy
runtime dependencies (NumPy, BLAS, LAPACK) or cover only Gaussian families.
We present \textbf{Gemmulem}, a pure-C library implementing Expectation--Maximisation
(EM) deconvolution for 35 univariate distribution families and three classes
of multivariate Gaussian and Student-$t$ mixtures.
Gemmulem requires no external dependencies and achieves competitive or
superior wall-clock performance through SIMD-accelerated E-steps
(AVX2/SSE2/scalar auto-dispatch), optional GPU offloading via OpenCL
for large datasets ($n \geq 50{,}000$), SQUAREM convergence acceleration,
streaming EM for out-of-memory data, and a statistically robust
\mbox{k-means++} initialiser backed by an xorshift128+ PRNG.
In benchmarks against \texttt{sklearn.mixture.GaussianMixture} across eight
scenarios (50 seeds, 95\,\% confidence intervals), Gemmulem is faster in
every case (1.5$\times$--6.6$\times$) while matching sklearn accuracy in
7 of 8 scenarios.
Gemmulem is released under GPL\,v3 at
\url{https://github.com/mathornton01/gemmulem}.
\end{abstract}

% =============================================================================
\section{Introduction}

Finite mixture models (FMMs) represent a distribution as a convex combination
of component densities:
\begin{equation}
  p(x \mid \Theta) = \sum_{k=1}^{K} \pi_k \, f_k(x \mid \theta_k),
  \quad \sum_{k=1}^{K} \pi_k = 1,
  \label{eq:fmm}
\end{equation}
where $\pi_k$ are mixture weights and $f_k$ are parametric densities.
FMMs underpin applications ranging from flow cytometry and astronomical
spectroscopy to speaker diarisation and genomics.

The Expectation--Maximisation algorithm of \citet{dempster1977maximum}
provides the standard approach to maximum-likelihood estimation of $\Theta$.
Despite widespread use, practical software packages exhibit two common
limitations: (i) restriction to the Gaussian family alone, and (ii)
dependence on large numerical-computing stacks (BLAS/LAPACK or NumPy/SciPy)
that complicate deployment in embedded systems, containerised environments,
and systems-level pipelines.

Python's \texttt{scikit-learn}~\citep{pedregosa2011scikit} offers a polished
Gaussian mixture implementation but covers no other distribution families.
Specialised R packages (\texttt{mixtools}, \texttt{flexmix}) widen family
coverage but require the R runtime and do not expose a C~API.
Proprietary solutions (MATLAB Statistics Toolbox, SAS PROC MIXED) add further
licensing constraints.

We introduce \textbf{Gemmulem}, a BSD-style C~library (C99) that addresses
both limitations: it supports 35 univariate families and full
multivariate Gaussian/Student-$t$ mixtures, compiles with a single
\texttt{make} invocation, links zero external libraries, and employs
several implementation techniques that confer substantial speed advantages
in practice.

% =============================================================================
\section{Implementation}

\subsection{Architecture}

Gemmulem is organised as a single-library compilation unit
(\texttt{libgemmulem.a} / \texttt{libgemmulem.so}) with a thin public header
(\texttt{gemmulem.h}).
The library is written in ISO C99; optional SIMD code paths use compiler
intrinsics guarded by \texttt{\#ifdef} macros so that the fallback scalar
path is always available.
The top-level API follows a \emph{context} pattern: callers allocate a
\texttt{gml\_ctx\_t}, populate it with data and hyperparameters, call
\texttt{gml\_fit()}, and inspect the resulting \texttt{gml\_result\_t}.

\subsection{Initialisation: k-means++ with Spectral Seeding}

Poor initialisation is a principal cause of convergence to local optima in
EM.
Gemmulem uses the k-means++ seeding strategy~\citep{arthur2007kmeans},
which selects initial centres with probability proportional to the squared
distance from already-selected centres, guaranteeing an $O(\log K)$ bound on
the initial cost relative to the global optimum in expectation.
Five independent restarts are performed; the run with the highest
log-likelihood is retained.

Random numbers are drawn from an xorshift128+ PRNG, which passes all
BigCrush statistical tests~\citep{lecuyer2007testu01} while requiring only
two 64-bit words of state and executing in two integer operations.

For univariate data, an additional \emph{spectral initialisation} pass
performs gap-detection on the sorted sample, providing a data-driven prior on
component count that reduces restart failures in well-separated mixtures.

\subsection{SIMD-Accelerated E-Step}

The E-step---computing responsibility weights
$r_{ik} = \pi_k f_k(x_i) / \sum_j \pi_j f_j(x_i)$---is the innermost loop
and dominates runtime for large $n$.
Gemmulem auto-dispatches to one of three code paths at library load time
using CPU feature detection:

\begin{enumerate}[leftmargin=*,noitemsep]
  \item \textbf{AVX2}: 256-bit SIMD, processes 8 doubles per cycle with
        cache-tiled blocking (128-row tiles sized for L1 data cache).
  \item \textbf{SSE2}: 128-bit fallback, processes 4 doubles per cycle.
  \item \textbf{Scalar}: ISO C, no intrinsics.
\end{enumerate}

Cache tiling ensures that the $n \times K$ responsibility matrix is traversed
in L1-friendly blocks, reducing cache-miss penalties for datasets with
$K > 8$ components.

\subsection{GPU Offloading via OpenCL}

For datasets with $n \geq 50{,}000$ observations, Gemmulem optionally
offloads the E-step to the first available OpenCL device (AMD or NVIDIA GPU,
or CPU fallback).
At library initialisation, the runtime probes for a platform; if none is
found the SIMD path is used transparently.
The OpenCL kernel replicates the cache-tiled algorithm in global memory with
local-memory reduction, achieving near-linear scaling with $n$ on supported
hardware.

\subsection{SQUAREM Convergence Acceleration}

Standard EM converges at a linear rate governed by the \emph{fraction of
missing information}.
For slow-converging families (e.g., highly overlapping Gaussians),
Gemmulem activates Squared Iterative Methods
(SQUAREM)~\citep{varadhan2008squarem} after iteration~30.
SQUAREM replaces one EM step with a quadratic extrapolation step
$\theta^{(t+1)} = \theta^{(t)} - 2\alpha \mathbf{r}_1 + \alpha^2 \mathbf{r}_2$,
where $\mathbf{r}_1, \mathbf{r}_2$ are successive EM increments and
$\alpha$ is chosen by an Armijo line-search.
This achieves super-linear (quasi-Newton) convergence without requiring
second-order derivatives.

\subsection{Streaming EM}

For datasets that exceed available RAM, Gemmulem provides a \emph{streaming
EM} interface: a callback supplies fixed-size data chunks from a file or
network stream, and the sufficient statistics for the M-step are accumulated
across chunks.
This extends Gemmulem to arbitrarily large datasets at the cost of
additional I/O passes per iteration.

% =============================================================================
\section{Features}

\subsection{Distribution Families}

Table~\ref{tab:families} lists the 35 univariate distribution families
currently supported.
Notably, Pearson system distributions (Types~I--VII) are included with
automatic type classification based on skewness and excess kurtosis of
component estimates.

\begin{table}[H]
\centering
\caption{Univariate distribution families supported by Gemmulem.}
\label{tab:families}
\small
\begin{tabular}{ll}
\toprule
\textbf{Family} & \textbf{Parameters} \\
\midrule
Gaussian (Normal)       & $\mu, \sigma^2$ \\
Exponential             & $\lambda$ \\
Gamma                   & $\alpha, \beta$ \\
Beta                    & $\alpha, \beta$ \\
Poisson                 & $\lambda$ \\
Weibull                 & $k, \lambda$ \\
Log-Normal              & $\mu, \sigma^2$ \\
Student-$t$             & $\nu, \mu, \sigma$ \\
Cauchy                  & $x_0, \gamma$ \\
Pearson I (Beta)        & auto-classified \\
Pearson II (sym.\ Beta) & auto-classified \\
Pearson III (Gamma)     & auto-classified \\
Pearson IV              & auto-classified \\
Pearson V (Inv.-Gamma)  & auto-classified \\
Pearson VI (Beta-Prime) & auto-classified \\
Pearson VII (gen.\ $t$) & auto-classified \\
Laplace                 & $\mu, b$ \\
Logistic                & $\mu, s$ \\
Uniform                 & $a, b$ \\
Pareto                  & $x_m, \alpha$ \\
Chi-Squared             & $k$ \\
Inverse Gaussian        & $\mu, \lambda$ \\
Rayleigh                & $\sigma$ \\
Maxwell-Boltzmann       & $a$ \\
Gompertz                & $\eta, b$ \\
Gumbel                  & $\mu, \beta$ \\
Fr\'{e}chet             & $\alpha, s, m$ \\
Burr XII                & $c, k$ \\
Dagum                   & $p, a, b$ \\
Benini                  & $\alpha, \beta, \sigma$ \\
Kumaraswamy             & $a, b$ \\
Power-Law               & $x_{\min}, \alpha$ \\
Half-Normal             & $\sigma$ \\
Triangular              & $a, b, c$ \\
von Mises               & $\mu, \kappa$ \\
\bottomrule
\end{tabular}
\end{table}

\subsection{Multivariate Mixtures}

For multivariate data, Gemmulem supports:
\begin{itemize}[leftmargin=*,noitemsep]
  \item \textbf{Multivariate Gaussian} with full, diagonal, and spherical
        covariance structures.
  \item \textbf{Multivariate Student-$t$} with shared or per-component
        degrees of freedom.
\end{itemize}

Covariance matrices are updated via Cholesky decomposition with a diagonal
regularisation term $\epsilon \mathbf{I}$ to ensure positive definiteness.

\subsection{Model Selection}

Selecting the number of components $K$ is a persistent challenge in mixture
modelling.
Gemmulem supports five criteria evaluated over a user-specified
range $k = 1, \ldots, K_{\max}$:

\begin{enumerate}[leftmargin=*,noitemsep]
  \item \textbf{BIC}~\citep{schwarz1978estimating}: $-2\ell + p\ln n$
  \item \textbf{AIC}: $-2\ell + 2p$
  \item \textbf{ICL} (Integrated Complete-data Likelihood):
        BIC penalised by component entropy
  \item \textbf{VBEM} (Variational Bayes EM): evidence lower-bound
  \item \textbf{MML} (Minimum Message Length): Wallace--Freeman coding cost
\end{enumerate}

The BIC-driven scan is the default for multivariate fits; the caller may
override the criterion via \texttt{gml\_ctx\_t.model\_sel}.

% =============================================================================
\section{Benchmarks}

\subsection{Experimental Setup}

We compare Gemmulem against
\texttt{sklearn.mixture.GaussianMixture}~\citep{pedregosa2011scikit}
on eight synthetic Gaussian mixture scenarios spanning a range of practical
conditions.
Each scenario was run over 50 independent random seeds; 95\,\% confidence
intervals (CIs) are constructed by the percentile bootstrap method.
Timing was measured as total wall-clock time from data ingestion to
convergence on a single CPU thread (Intel i9-10850K, AVX2 enabled;
Gemmulem SIMD path active, OpenCL disabled for fair comparison).
Accuracy is quantified as the mean absolute error (MAE) between
true and estimated means $\bar{\delta}_\mu$, variances $\bar{\delta}_\sigma$,
and mixture weights $\bar{\delta}_\pi$, after optimal component matching
by the Hungarian algorithm.

\subsection{Results}

Table~\ref{tab:benchmarks} summarises the results.
Gemmulem achieves strictly faster runtime in all eight scenarios (1.5$\times$
to 6.6$\times$), with the largest gains for overlapping components where
SQUAREM acceleration is most impactful.
In 7 of 8 scenarios the 95\,\% CIs on all three MAE metrics overlap between
Gemmulem and sklearn, indicating statistically equivalent accuracy.
The single exception is the high-$K$, moderate-separation scenario
($K = 8$, $5\sigma$), where sklearn shows a marginal edge on variance
and weight MAE; this is attributable to sklearn's covariance regularisation
heuristic, which is tuned for the Gaussian case.

\begin{table*}[t]
\centering
\caption{Benchmark results: Gemmulem vs.\ \texttt{sklearn.mixture.GaussianMixture}.
  Speed is wall-clock speedup ratio (Gemmulem/sklearn$^{-1}$; higher is better for Gemmulem).
  Accuracy columns show mean absolute error with 95\,\% CI overlap (\checkmark = overlapping;
  $\dagger$ = sklearn slightly lower). $n = 10{,}000$ unless noted.}
\label{tab:benchmarks}
\small
\begin{tabular}{lccccc}
\toprule
\textbf{Scenario} & \textbf{$K$} & \textbf{Speedup} & \textbf{MAE $\mu$} & \textbf{MAE $\sigma^2$} & \textbf{MAE $\pi$} \\
\midrule
Well-separated ($8\sigma$)     & 3 & 3.7$\times$ & \checkmark & \checkmark & \checkmark \\
Overlapping ($2\sigma$)        & 3 & 6.6$\times$ & \checkmark & \checkmark & \checkmark \\
Unequal weights (70/20/10\,\%) & 3 & 4.7$\times$ & \checkmark & \checkmark & \checkmark \\
High-$K$ ($5\sigma$)           & 8 & 1.9$\times$ & \checkmark & $\dagger$  & $\dagger$  \\
Small sample ($n = 200$)       & 3 & 1.5$\times$ & \checkmark & \checkmark & \checkmark \\
Large sample ($n = 50{,}000$)  & 3 & 3.0$\times$ & \checkmark & \checkmark & \checkmark \\
Unequal variance               & 3 & 4.5$\times$ & \checkmark & \checkmark & \checkmark \\
High-$K$ overlapping ($3\sigma$) & 5 & 4.4$\times$ & \checkmark & \checkmark & \checkmark \\
\midrule
\multicolumn{2}{l}{\textit{Summary}} & \textbf{8/8} faster & \multicolumn{3}{c}{\textbf{7/8} scenarios fully tied} \\
\bottomrule
\end{tabular}
\end{table*}

The speedup for overlapping mixtures ($K=3$, $2\sigma$, $6.6\times$) is
primarily attributable to SQUAREM acceleration: these problems require
significantly more EM iterations before convergence, and SQUAREM's
super-linear convergence dramatically reduces the iteration count.
The more modest gain for small samples ($n=200$, $1.5\times$) reflects the
fact that at low $n$ the E-step itself is cheap; framework overhead
in sklearn becomes relatively less significant.

% =============================================================================
\section{Availability}

\subsection{Source Code and License}

Gemmulem is available under the GNU General Public License v3.0:
\begin{center}
  \url{https://github.com/mathornton01/gemmulem}
\end{center}

\subsection{Installation}

Building from source requires only a C99-compliant compiler
(GCC $\geq$ 4.9 or Clang $\geq$ 3.5) and GNU Make:

\begin{verbatim}
  git clone https://github.com/mathornton01/gemmulem
  cd gemmulem
  make          # builds libgemmulem.a and examples
  make install  # installs to /usr/local by default
\end{verbatim}

SIMD acceleration is enabled automatically when the target CPU supports
AVX2 or SSE2; no flags are required.
OpenCL support is compiled in when \texttt{libOpenCL} is present on the
include path; it is silently omitted otherwise.

\subsection{Usage Example}

The following minimal C snippet fits a 3-component Gaussian mixture:

\begin{verbatim}
  #include <gemmulem.h>

  gml_ctx_t ctx = gml_ctx_default();
  ctx.n = n;  ctx.data = x;
  ctx.K = 3;  ctx.family = GML_GAUSSIAN;
  ctx.max_iter = 500;  ctx.tol = 1e-6;

  gml_result_t res = gml_fit(&ctx);
  for (int k = 0; k < res.K; k++)
      printf("pi=%f mu=%f sigma2=%f\n",
             res.weights[k],
             res.params[k][0],
             res.params[k][1]);
  gml_result_free(&res);
\end{verbatim}

Language bindings for Python and R are under active development.

% =============================================================================
\section*{Acknowledgements}

The authors thank the contributors to the xorshift128+ PRNG literature and
the OpenCL working group whose open specifications made the GPU backend
feasible without proprietary SDK dependencies.

% =============================================================================
\bibliographystyle{plain}
\bibliography{refs}

\begin{thebibliography}{9}

\bibitem{dempster1977maximum}
A.~P. Dempster, N.~M. Laird, and D.~B. Rubin.
\newblock Maximum likelihood from incomplete data via the {EM} algorithm.
\newblock \emph{Journal of the Royal Statistical Society: Series B}, 39(1):1--38, 1977.

\bibitem{arthur2007kmeans}
D.~Arthur and S.~Vassilvitskii.
\newblock {k-means++}: The advantages of careful seeding.
\newblock In \emph{Proceedings of the Eighteenth Annual ACM-SIAM Symposium on
  Discrete Algorithms (SODA)}, pages 1027--1035, 2007.

\bibitem{varadhan2008squarem}
R.~Varadhan and C.~Roland.
\newblock Simple and globally convergent methods for accelerating the
  convergence of any {EM} algorithm.
\newblock \emph{Scandinavian Journal of Statistics}, 35(2):335--353, 2008.

\bibitem{schwarz1978estimating}
G.~Schwarz.
\newblock Estimating the dimension of a model.
\newblock \emph{The Annals of Statistics}, 6(2):461--464, 1978.

\bibitem{blackford2002updated}
L.~S. Blackford, J.~Demmel, J.~Dongarra, I.~Duff, S.~Hammarling,
  G.~Henry, M.~Heroux, L.~Kaufman, A.~Lumsdaine, A.~Petitet,
  R.~Pozo, K.~Remington, and R.~C. Whaley.
\newblock An updated set of basic linear algebra subprograms ({BLAS}).
\newblock \emph{ACM Transactions on Mathematical Software}, 28(2):135--151, 2002.

\bibitem{pedregosa2011scikit}
F.~Pedregosa, G.~Varoquaux, A.~Gramfort, V.~Michel, B.~Thirion, O.~Grisel,
  M.~Blondel, P.~Prettenhofer, R.~Weiss, V.~Dubourg, J.~Vanderplas,
  A.~Passos, D.~Cournapeau, M.~Brucher, M.~Perrot, and E.~Duchesnay.
\newblock Scikit-learn: Machine learning in {Python}.
\newblock \emph{Journal of Machine Learning Research}, 12:2825--2830, 2011.

\bibitem{lecuyer2007testu01}
P.~L'Ecuyer and R.~Simard.
\newblock {TestU01}: A {C} library for empirical testing of random number generators.
\newblock \emph{ACM Transactions on Mathematical Software}, 33(4):22, 2007.

\end{thebibliography}

\end{document}
